\section{导数的基础概念}

\subsection{平均变化率与瞬时变化率}

\begin{theorem}[平均变化率]
    函数 $f(x)$ 在区间 $[a, b]$ 上的平均变化率为 $\frac{\Delta y}{\Delta x} = \frac{f(b) - f(a)}{b - a}$.（即为割线的斜率）
\end{theorem}
\begin{example}
    函数 $f(x) = x^{2} - \cos x$ 在区间 $\left[ 0, \pi \right]$ 上的平均变化率为（\quad）
    \begin{tasks}(4)
        \task $- \pi - \frac{2}{\pi}$
        \task $- \pi$
        \task $\pi$
        \task $\pi + \frac{2}{\pi}$
    \end{tasks}
\end{example}

\v{12em}

\begin{example}
    在一次跳水运动中，某运动员在运动过程中的重心相对于水面的高度 $h$（单位： m ）与起跳后的时间 $t$（单位： s ）存在函数关系 $h(t)=-4.9 t^2+2.8 t+11$ ，则下列说法正确的是（ \quad）
    \begin{itemize}
        \item 在 $0 \leq t \leq 0.2$ 这段时间里，运动员的平均速度 $\bar{v}=1.82 \mathrm{~m} / \mathrm{s}$
        \item 在运动过程中运动员的瞬时速度 $v(t)=-9.8 t+2.8$
        \item 在起跳到落水的过程中运动员的速度不可能为 0
        \item 第 0.2 s 时刻瞬时速度为 $0.84 \mathrm{~m} / \mathrm{s}$
    \end{itemize}
\end{example}

\v{12em}

\begin{theorem}[瞬时速度的物理意义]
    物体在 $t=t_0$ 时刻的瞬时速度是其位移函数 $y(t)$ 在该点的导数，即 $v(t_0) = y'(t_0)$.
\end{theorem}

\begin{example}
    某物体沿直线运动，位移 $y$（单位：$\mathrm{m}$）与时间 $t$（单位：$\mathrm{s}$）的关系为 $y(t) = 1 - t + t^{2}$，则该物体在 $t = 3\,\mathrm{s}$ 时的瞬时速度是（\quad）
    \begin{tasks}(4)
        \task $2\,\mathrm{m/s}$
        \task $3\,\mathrm{m/s}$
        \task $4\,\mathrm{m/s}$
        \task $5\,\mathrm{m/s}$
    \end{tasks}
\end{example}

\v{12em}

\subsection{脱衣大法求函数定义}

\begin{theorem}[脱衣大法]
    \textbf{导数的定义及其变式}:
    $f'(x_0) = \lim_{\Delta x \to 0}\frac{f(x_0 + \Delta x) - f(x_0)}{\Delta x}$.
    可以有变式 $k \cdot f'(x_0) = \lim_{\Delta x \to 0}\frac{f(x_0) - f(x_0 - k \Delta x)}{\Delta x}$.

    任何这些有变式的形式，都可以想办法凑出来，但采用一种简单的"脱衣大法"来解决问题。
\end{theorem}
\begin{example}
    已知函数 $f(x) = \frac{1}{3}x^{3} - x$，则 $\lim_{\Delta x \to 0}\frac{f(3) - f(3 - 2 \Delta x)}{\Delta x} =$ （\quad）
    \begin{tasks}(4)
        \task 2 \task 4 \task 8 \task 16
    \end{tasks}
\end{example}

\v{12em}

\begin{example}
    已知函数 $f(x)$ 在 $x = x_{0}$ 处可导，且 $\lim_{\Delta x \to 0}\frac{f(x_{0} - \Delta x) - f(x_{0})}{\Delta x} = 2025$，则 $f'(x_{0}) =$ （\quad）
    \begin{tasks}(4)
        \task $- 2025$ \task 0 \task 1 \task 2025
    \end{tasks}
\end{example}

\v{12em}

\begin{example}
    设 $f(x)$ 是可导函数，若 $\lim_{\Delta x \to 0}\frac{f(3 - 3\Delta x) - f(3)}{\Delta x} = 3$，则 $f'(3)$ （\quad）
    \begin{tasks}(4)
        \task $- 1$ \task $- \frac{1}{3}$ \task $\frac{1}{3}$ \task 1
    \end{tasks}
\end{example}

\v{12em}